\section*{Chapitre \ref{ch:hf:pnl-analytics-and-the-strategy-lifecycle} --- Analyse du P\&L et cycle de vie d'une stratégie}

\begin{solution}{exo:hf:pnl-analytics-and-the-strategy-lifecycle:1}
$15\,000\,000\times0.12\times(0.8\times1-0.35)/100=\$8\,100$ par jour.
\end{solution}

\begin{solution}{exo:hf:pnl-analytics-and-the-strategy-lifecycle:2}
$15\,000\,000\times0.084\times(0.68-0.35)/100=\$4\,158$ par jour.
\end{solution}

\begin{solution}{exo:hf:pnl-analytics-and-the-strategy-lifecycle:3}
$15\,000\,000\times0.132\times(0.80-0.45)/100=\$6\,930$, inférieur à 8\,100~\$~: il achète 10~\% de volume en plus, à une marge plus mince de 22~\%.
\end{solution}

\begin{solution}{exo:hf:pnl-analytics-and-the-strategy-lifecycle:4}
Tout le reste qui a changé après le douzième (fourchettes, volumes, un concurrent) modifie le P\&L dans le même temps~; avant et après ne permettent pas
de les séparer. Randomiser les jours (ou les instruments) qui subissent le changement permet de l'isoler de tout le reste.
\end{solution}

\begin{solution}{exo:hf:pnl-analytics-and-the-strategy-lifecycle:5}
Le concurrent réduit une part de la marge~; celle-ci est plus faible une fois la fourchette resserrée. Évalué après le changement de fourchette, l'effet du
concurrent est plus faible que s'il était évalué en premier.
\end{solution}

\begin{solution}{exo:hf:pnl-analytics-and-the-strategy-lifecycle:6}
Elles sont mesurées sur chaque exécution, avec beaucoup moins de bruit que le P\&L d'une journée, qui intègre l'inventaire et les mouvements de marché~; leur effet sur le P\&L
découle alors du produit.
\end{solution}

\begin{solution}{exo:hf:pnl-analytics-and-the-strategy-lifecycle:7}
Graine 2~: résidu 6\,400~\$ injecté (2\,600~\$ estimé)~; concurrent $-\$53\,400$ estimé contre $-\$57\,300$~; paramètre $-\$8\,100$ contre $-\$8\,000$.
Graine 3~: résidu $-\$3\,600$ ($-\$7\,600$)~; concurrent $-\$55\,400$ contre $-\$59\,400$~; paramètre $-\$7\,900$ contre $-\$7\,900$. Le concurrent est
sous-estimé d'environ 4\,000~\$ à chaque fois, le seul jour où le point de rupture arrive en retard.
\end{solution}

\begin{solution}{exo:hf:pnl-analytics-and-the-strategy-lifecycle:8}
Au mois~9 simulé, la fourchette s'est resserrée et un concurrent est arrivé dans les mêmes semaines~: ils ont coûté six fois plus que le paramètre. Avant et
après ne permettent pas de dire lequel en est responsable~; l'expérience et les mesures par facteur, elles, le permettent.
\end{solution}

\begin{solution}{pb:hf:pnl-analytics-and-the-strategy-lifecycle:1}
\begin{enumerate}
\item $V\,s\,(\kappa h-a)$~: volume de marché et demi-fourchette issus du marché, part issue des exécutions, ratio de capture à partir des prix d'exécution, sélection
adverse à partir des mark-outs.
\item Deux plateformes~; fourchette plus étroite et volume plus élevé au mois~9~; un concurrent sur la plateforme~B au jour~175~; un changement de paramètre sur des jours randomisés à partir du jour~182.
\item Modifier un facteur à la fois, du mois~8 au mois~9, dans un ordre fixe~; le résidu est ce qui reste.
\item Un teneur de marché a gagné 1.55~\euro par transaction sur la fourchette et perdu 0.68~\euro sur ses positions.
\item Voir la \cref{def:hf:pnl-analytics-and-the-strategy-lifecycle:control}.
\item $-\$2\,670$ par jour (erreur type 1\,070~\$)~; avec CUPED $-\$2\,550$ (1\,070~\$)~; part $+10~\%$ et sélection adverse $+0.10$ centime exactement.
\item Voir la \cref{def:hf:pnl-analytics-and-the-strategy-lifecycle:decay}~; un test CUSUM unilatéral sur la part de la plateforme~B franchit son seuil le jour~176.
\item 29.9~\% de part et 14.9~\% de ratio de capture, contre 30~\% et 15~\%.
\item Injectés~: fourchette $-\$51\,500$, volume $+\$48\,000$, concurrent $-\$57\,400$, paramètre $-\$8\,100$, résidu $+\$23\,400$. Estimés~: les mêmes
sauf le concurrent $-\$53\,500$ et le résidu $+\$19\,600$.
\item Le point de rupture est en retard d'un jour, donc l'effet du concurrent s'applique à un jour de moins.
\item Le bruit de l'inventaire, du volume et des fourchettes, et l'approximation des moyennes mensuelles.
\item L'attribution en dépend~; un rapport doit être comparable d'un mois à l'autre.
\item Sur la baisse de 45\,600~\$ du mois~9~: fourchette $-\$51\,500$, volume $+\$48\,000$, concurrence $-\$53\,500$ estimée ($-\$57\,400$ injectée), paramètre
$-\$8\,100$ (les deux), résidu d'environ 20\,000~\$.
\item Non~: il coûte 2\,670~\$ par jour.
\item Voir la \cref{def:hf:pnl-analytics-and-the-strategy-lifecycle:lifecycle}.
\item Jour 239 pour un coût de 12\,000~\$ par jour, avec une fenêtre de 60 jours~; cela dépend du coût et de la fenêtre.
\item Le P\&L journalier par facteur, par plateforme et par instrument~; les points de rupture~; les expériences en cours~; le résidu.
\item Elle moyenne la contribution de chaque facteur sur tous les ordres possibles, ce qui supprime l'arbitraire de l'ordre des facteurs.
\item Ses modifications se confondraient avec l'arrivée du concurrent.
\item Savoir sur quelle cause agir.
\end{enumerate}
\end{solution}

\begin{solution}{iq:hf:pnl-analytics-and-the-strategy-lifecycle:1}
Les facteurs~: volume de marché et fourchettes, part et ratio de capture par plateforme, sélection adverse par mark-out, P\&L d'inventaire, frais~; puis toute modification
apportée à la stratégie.
\iqlookfor{décomposer avant d'expliquer.}
\end{solution}

\begin{solution}{iq:hf:pnl-analytics-and-the-strategy-lifecycle:2}
Un point de rupture dans la part et le ratio de capture sur cette seule plateforme, avec un volume de marché inchangé, et des positions dans la file qui commencent plus tard.
\iqlookfor{points de rupture propres à une plateforme.}
\end{solution}

\begin{solution}{iq:hf:pnl-analytics-and-the-strategy-lifecycle:3}
Un journal des changements indiquant qui, quoi, quand, pourquoi et l'assignation de l'expérience~; des traitements nocturnes rapprochant exécutions, mark-outs, données de marché et assignations~; le
rapport construit à partir d'eux, avec le résidu.
\iqlookfor{piste d'audit et rapports reproductibles.}
\end{solution}

\begin{solution}{iq:hf:pnl-analytics-and-the-strategy-lifecycle:4}
Un motif écrit, une revue, un moyen de revenir en arrière, et une randomisation ou un périmètre limité pour pouvoir mesurer son effet~; les exceptions ne sont admises que pour réduire le risque.
\iqlookfor{contrôle des changements.}
\end{solution}

\begin{solution}{iq:hf:pnl-analytics-and-the-strategy-lifecycle:5}
Environ $2(1.96+0.84)^2\times4\,500^2/2\,000^2\approx80$ jours par bras, 159 au total.
\iqlookfor{la formule de taille d'échantillon.}
\end{solution}

\begin{solution}{iq:hf:pnl-analytics-and-the-strategy-lifecycle:6}
$\Delta(xy)=\Delta x\,y_0+x_1\Delta y$~: changer $x$ d'abord lui attribue $\Delta x\,y_0$ et $y$ reçoit $x_1\Delta y=x_0\Delta y+\Delta x\Delta y$, le
terme d'interaction~; l'ordre inverse l'attribue à $x$.
\iqlookfor{le terme d'interaction.}
\end{solution}
